\section*{Chapter \ref{ch:g11:plate-motions} --- Plates and Their Motions}

\begin{solution}{exo:g11:plate-motions:1}
A \omterm{def:g11:plate-motions:rigid-plate}{rigid plate} is a piece of \omterm{def:g11:lithosphere-asthenosphere:lithosphere}{lithosphere} that moves as a single block: the
distances between its points do not change, and its deformation is gathered in
narrow zones along its boundaries. The earthquakes show it: almost all of them
happen along narrow belts, the plate boundaries, and very few inside the
plates.
\end{solution}

\begin{solution}{exo:g11:plate-motions:2}
\qty{72}{mm} a year is \qty{72}{km} per million years. In 30 years the station
moves $72 \times 30 = \qty{2160}{mm}$, about \qty{2.2}{m}; in a million years,
\qty{72}{km}.
\end{solution}

\begin{solution}{exo:g11:plate-motions:3}
$v_E = 30/10 = \qty{3}{cm}$ a year and $v_N = 40/10 = \qty{4}{cm}$ a year. The
speed is $\sqrt{3^2 + 4^2} = \qty{5}{cm}$ a year. The angle $\beta$ with the
north satisfies $\tan\beta = 3/4$, so $\beta \approx \ang{37}$: the station
moves toward the north-east, \ang{37} east of north.
\end{solution}

\begin{solution}{exo:g11:plate-motions:4}
A \omterm{def:g11:plate-motions:hotspot-track}{hotspot track} is the line of volcanoes and seamounts that a plate carries
away from a hotspot as it moves over it, older and older with distance from
the active volcano. The youngest volcano, Kilauea, stands at the
south-eastern end of the chain, on the island of Hawaii; the ages increase
toward the west-north-west. The plate carries its volcanoes away from the
hotspot in the direction of its motion: the Pacific plate moves toward the
west-north-west.
\end{solution}

\begin{solution}{exo:g11:plate-motions:5}
$519/5.1 \approx \qty{102}{km}$ per million years, that is \qty{102}{mm} a
year, about \qty{10}{cm} a year.
\end{solution}

\begin{solution}{exo:g11:plate-motions:6}
$v_E = 63/31 \approx \qty{20.3}{mm}$ a year and $v_N = 48/31 \approx
\qty{15.5}{mm}$ a year. The speed is $\sqrt{20.3^2 + 15.5^2} \approx
\qty{25.5}{mm}$ a year; $\tan\alpha = 15.5/20.3 \approx 0.76$ gives $\alpha
\approx \ang{37}$: the station moves \ang{37} north of east, toward the
north-east. Mauna Kea moves $72/25.5 \approx 2.8$ times faster.
\end{solution}

\begin{solution}{exo:g11:plate-motions:7}
(1) In millimetres a year,
\[
  \vec v_{\mathrm{NA/EU}} = (-10.5 - 7.7,\ 20.6 - 16.3) = (-18.2, +4.3).
\] Its speed is $\sqrt{18.2^2 + 4.3^2} \approx
\qty{18.7}{mm}$ a year, and $\tan\alpha = 4.3/18.2 \approx 0.24$ gives
\ang{13} north of west. (2) The two plates move apart across the ridge at
about \qty{1.9}{cm} a year: this is the \omterm{def:g11:drift-to-spreading:spreading-rate}{full spreading rate} of the ridge in
Iceland. (3) It agrees with the \qty{2}{cm} a year quoted for the North
Atlantic.
\end{solution}

\begin{solution}{exo:g11:plate-motions:8}
(1) $v = \pi \times 0.651 \times 6371/180 \approx \qty{72.4}{km}$ per million
years, about \qty{7.2}{cm} a year. (2) $72.4 \times \sin\ang{30} \approx
\qty{36.2}{km}$ per million years, about \qty{3.6}{cm} a year. (3) At the
rotation pole, if it lies on the plate. The speed changes from point to
point, from zero at the pole to its largest value \ang{90} away from it: a
whole plate has no single speed, only a speed at each point (or an average
over its surface).
\end{solution}

\begin{solution}{exo:g11:plate-motions:9}
(1) In millimetres a year,
\[
  \vec v_{\mathrm{IN/EU}} = (40.6 - 24.0,\ 38.6 + 3.3) = (+16.6, +41.9).
\] The speed is $\sqrt{16.6^2 + 41.9^2} \approx
\qty{45}{mm}$ a year; the angle $\beta$ with the north satisfies $\tan\beta =
16.6/41.9 \approx 0.40$, so $\beta \approx \ang{22}$: India moves toward
Eurasia, \ang{22} east of north. (2) \qty{45}{km} per million years, so
\qty{450}{km} in 10 million years.
\end{solution}

\begin{solution}{exo:g11:plate-motions:10}
$(2432 - 1058)/(27.7 - 10.3) = 1374/17.4 \approx \qty{79}{km}$ per million
years, about \qty{7.9}{cm} a year, a little less than \qty{9.1}{cm} a year.
The ages are uncertain by about half a million years each; the age
difference, $17.4 \pm 1.0$ million years, gives a speed between about 75 and
\qty{84}{km} per million years. Between these two volcanoes the plate seems
to have moved a little more slowly than on average.
\end{solution}

\begin{solution}{exo:g11:plate-motions:11}
(1) $\vec v_{A/C} = \vec v_{A/B} + \vec v_{B/C}$: \qty{3}{cm} a year east plus
\qty{4}{cm} a year north. The two arrows form a right angle, and the third
side of the triangle has the length $\sqrt{3^2 + 4^2} = \qty{5}{cm}$ a year,
in the direction \ang{37} east of north ($\tan\beta = 3/4$). (2) $\vec v_{C/A}
= -\vec v_{A/C}$: \qty{5}{cm} a year in the opposite direction, \ang{37} west
of south.
\end{solution}

\begin{solution}{exo:g11:plate-motions:12}
(1) $54/0.375 = \qty{144}{km}$ per million years, \qty{14.4}{cm} a year. (2)
From $29/0.375 \approx 77$ to $79/0.375 \approx \qty{211}{km}$ per million
years: between about 8 and \qty{21}{cm} a year. (3) Close to the hotspot the
distance is no larger than the size of a volcano or of the region over which
the hotspot feeds volcanoes, and the age is no larger than the time a volcano
takes to grow: both are very uncertain in relative terms. Far from the
hotspot an error of \qty{25}{km} on \qty{2400}{km}, or of half a million years
on 28 million, is only one or two per cent.
\end{solution}

\begin{solution}{exo:g11:plate-motions:13}
(1) Ten times the scatter is \qty{50}{mm}: $50/72 \approx 0.7$ years, about
eight months. (2) $50/5 = 10$ years. (3) Over millions of years the plates move
tens to hundreds of kilometres, distances that the magnetic stripes record and
that were easy to measure at sea; positioning precise to a few millimetres,
which makes a motion visible in months, came much later.
\end{solution}

\begin{solution}{exo:g11:plate-motions:14}
(1) In millimetres a year,
\[
  \vec v_{\mathrm{SA/AF}} = (-4.9 - 27.1,\ 11.0 - 17.9) = (-32.0, -6.9).
\] The speed is $\sqrt{32.0^2 + 6.9^2} \approx
\qty{32.7}{mm}$ a year, and $\tan\alpha = 6.9/32.0 \approx 0.22$ gives
\ang{12} south of west. (2) The \omterm{def:g11:drift-to-spreading:spreading-rate}{full spreading rate} is about \qty{3.3}{cm} a
year, the half rate about \qty{1.6}{cm} a year. (3) The half rate of the last
few million years, \qty{1.6}{cm} a year, is a little smaller than the
\qty{1.9}{cm} a year averaged over 66 million years: the South Atlantic may
have opened slightly more slowly recently than on average, or the difference
may come from the uncertainties of the sediment ages. Both measurements agree
that the ridge is slow.
\end{solution}

\begin{solution}{exo:g11:plate-motions:15}
(1) A straight line rising by \qty{20}{mm} a year, a sudden drop of
\qty{50}{mm} after five years, then a line rising again with the same slope.
(2) The station ends at $10 \times 20 - 50 = \qty{150}{mm}$ east of its start:
$150/10 = \qty{15}{mm}$ a year, a quarter too small. (3) Each half gives
\qty{20}{mm} a year. (4) The jump is a separate, sudden motion caused by the
earthquake: the geodesist measures its size from the positions just before and
just after it, removes it from the series (or fits the two parts with the
same slope and a step between them), and keeps \qty{20}{mm} a year as the
velocity of the station.
\end{solution}

\begin{solution}{pb:g11:plate-motions:1}
\textbf{1.} Mauna Kea (\num{0.375} million years) is the youngest, Suiko
(centre) (\num{64.7} million years) the oldest. The ages increase with the
distance from Kilauea.

\textbf{2.} $\qty{1}{km} = 10^6~\unit{mm}$ and one million years is $10^6$
years: \qty{1}{km} per million years is $10^6~\unit{mm}/10^6$ years,
\qty{1}{mm} a year.

\textbf{3.} $519/5.1 \approx \qty{102}{mm}$ a year.

\textbf{4.} Necker $1058/10.3 \approx 102.7$; Laysan $1818/19.9 \approx
91.4$; Midway $2432/27.7 \approx 87.8$; Colahan $3128/38.6 \approx 81.0$;
Abbott $3280/38.7 \approx 84.8$ (millimetres a year).

\textbf{5.} The points (54; 0.375), (519; 5.1), (780; 7.2), (1058; 10.3),
(1818; 19.9), (2432; 27.7), (3128; 38.6), (3280; 38.7) lie close to a straight
line through the origin, with some scatter (against a line of \qty{91}{km} per
million years, Kauai, Nihoa and Necker lie a little below it, Laysan on it,
Colahan and Abbott a little above).

\textbf{6.} Each volcano was built over the hotspot and has moved with the
plate since then. In the time $t$, the plate has moved $v\,t$ in a steady
direction, so the volcano lies at $d = v\,t$ from the point now over the
hotspot: $t = d/v$.

\textbf{7.} $(102.7 + 91.4 + 87.8 + 81.0 + 84.8)/5 \approx \qty{89.5}{mm}$ a
year.

\textbf{8.} 89.5 against \qty{91}{km} per million years: they differ
by less than \qty{2}{\percent}.

\textbf{9.} $519/91 \approx 5.7$ million years, against an age of \num{5.1}
million years: close, Kauai being slightly younger than the average line
predicts.

\textbf{10.} The volcanoes are older and older toward the west-north-west:
the plate carries them in that direction. The Pacific plate moves toward the
west-north-west over the hotspot.

\textbf{11.} $(4860 - 3493)/(64.7 - 42.4) = 1367/22.3 \approx \qty{61}{km}$
per million years, about \qty{6.1}{cm} a year.

\textbf{12.} $65/91 \approx 0.71$: along the Emperor chain the plate moved
about \qty{30}{\percent} more slowly than along the Hawaiian chain.

\textbf{13.} $3668/50 \approx \qty{73}{km}$ per million years, about
\qty{7.3}{cm} a year, less than \qty{91}{km} per million years. The line through the origin
uses the 1987 ages, which put the end of the Hawaiian chain at 39 to 43
million years; the newer age of about 50 million years for the start of the
bend spreads the same distance over a longer time. Both are right if the plate
moved more slowly in the early part of the Hawaiian chain and faster later:
the speed was not constant.

\textbf{14.} A change in the direction of motion of the Pacific plate, about
50 million years ago.

\textbf{15.} The direction of the chain is the direction in which the plate
moved; the spacing of the ages along it is its speed. A plate can speed up or
slow down without turning, or turn without changing speed, so the two are
independent. Changes of speed with no change of direction along the Emperor
chain show that its speed record and the bend are separate events: the bend
was not made by the hotspot slowing down.

\textbf{16.} $\sqrt{63.2^2 + 34.5^2} \approx \qty{72.0}{mm}$ a year, and
$\tan\alpha = 34.5/63.2 \approx 0.55$ gives \ang{29} north of west.

\textbf{17.} West-north-west is \ang{22.5} north of west, north-west
\ang{45}: the station's direction, \ang{29} north of west, lies between the
two, close to the direction of the chain. It is consistent.

\textbf{18.} The station's velocity is measured with respect to a frame in
which the \omterm{def:g11:lithosphere-asthenosphere:lithosphere}{lithosphere} as a whole does not turn, while the chain gives the
motion with respect to the hotspot: if the two frames move with respect to
each other, the speeds differ. And the station measures the last 30 years,
the chain the last 40 million years, during which the speed may have changed.
The scatter of the ages also leaves an uncertainty of a few millimetres a
year on the speed of the chain.

\textbf{19.} $100/91 \approx 1.1$ million years.

\textbf{20.} \textbf{The Pacific plate has moved over the Hawaiian hotspot at
about \qty{9}{cm} a year} during the last 40 million years, toward the
west-north-west; before the bend it moved more slowly, and in another
direction.
\end{solution}
